\chapter{常见错误速查}
\label{ch:appA}

诊断先按发生阶段定位：解析检查语法结构，语义检查类型与借用，链接阶段查找目标库和符号，运行阶段检查实际输入与资源。下面保留可执行的拒绝示例及其关键诊断；文件名、位置和调用栈会随保存位置变化。

\section{定位阶段}

\begin{center}
\begin{tabular}{lll}
\toprule
阶段 & 检查内容 & 后续动作 \\
\midrule
解析 & 标点、括号与声明结构 & 从首个语法错误修正上下文 \\
类型与借用 & 实参与返回类型、引用来源、资源使用 & 检查声明与存活范围 \\
代码生成与链接 & 目标架构、ABI、库与符号 & 核对工具链和真实依赖 \\
应用运行 & 输入、边界、缺省与外部资源 & 建立前提或返回可恢复错误 \\
\bottomrule
\end{tabular}
\end{center}

\section{除零与运行期前提}

先检查分母再执行除法。输入解析成功并不保证它可以作为分母。

\par
\medskip
\begin{lstlisting}[style=chali]
void main()
{
    int64 a = 1;
    int64 b = 0;
    println("{}", a / b);
}
\end{lstlisting}
\par
\medskip

\par
\medskip
\begin{lstlisting}[style=out]
chali aot: 整数除零
出错位置：hello.chl:5（main）
调用栈（内 → 外）：
  main
\end{lstlisting}
\par
\medskip

\section{缩窄转换}

确认目标范围，再显式转换；不能仅为消除诊断而截断业务数据。

\par
\medskip
\begin{lstlisting}[style=chali]
void main()
{
    int64 big = 70000;
    int32 small = big;
    // 缩窄：位宽变小
    println("{}", small);
}
\end{lstlisting}
\par
\medskip

\par
\medskip
\begin{lstlisting}[style=out]
chali: hello.chl 第 4 行第 5 列：不能隐式转换：int64 -> int32（请显式强转）
\end{lstlisting}
\par
\medskip

\section{字面量越界}

选择能容纳数值的类型，或修正本来就不合理的常量。

\par
\medskip
\begin{lstlisting}[style=chali]
void main()
{
    uint8 b = 256;
    // uint8 只到 255
}
\end{lstlisting}
\par
\medskip

\par
\medskip
\begin{lstlisting}[style=out]
chali: hello.chl 第 3 行第 5 列：整数字面量 256 超出 uint8 的取值范围（0..255）
\end{lstlisting}
\par
\medskip

\section{引用来源}

引用必须绑定允许借用的实际对象，数值转换不能制造合法引用。

\par
\medskip
\begin{lstlisting}[style=chali]
void main()
{
    int64 mut& r = 5;
    // 引用要从对象身上取地址，不能绑字面量
}
\end{lstlisting}
\par
\medskip

\par
\medskip
\begin{lstlisting}[style=out]
chali: hello.chl 第 3 行第 5 列：不能隐式转换：int64 -> int64 mut&（请显式强转）
\end{lstlisting}
\par
\medskip

\section{未绑定引用}

在声明处绑定对象，并确保其存活覆盖引用的使用。

\par
\medskip
\begin{lstlisting}[style=chali]
void main()
{
    int64 mut& r;
    println("{}", r);
}
\end{lstlisting}
\par
\medskip

\par
\medskip
\begin{lstlisting}[style=out]
chali: hello.chl 第 3 行第 5 列：引用声明必须初始化，缺省值请使用 Option<T mut&>
\end{lstlisting}
\par
\medskip

\section{缺省值取出}

先匹配 Option 的状态或检查 isNone，再使用取值接口。需要缺省时显式构造 None；默认构造不应被当作缺省。

\par
\medskip
\begin{lstlisting}[style=chali]
void main()
{
    Option<int64> none = None;
    println("{}", none.unwrap());
}
\end{lstlisting}
\par
\medskip

\par
\medskip
\begin{lstlisting}[style=out]
chali aot: panic：called Option.unwrap() on None
\end{lstlisting}
\par
\medskip

\section{排查顺序}

先保存触发问题的最小程序，保留编译器版本、命令、配置和首条诊断。把外部输入、文件、网络或数据库问题与语言检查区分开。改动后重新验证原失败条件和正常路径；未运行的检查不写成通过。更多诊断和接口边界见\href{https://github.com/rabbitc999/ChaLi/blob/master/doc/04-language-reference.md}{语言参考}、\href{https://github.com/rabbitc999/ChaLi/blob/master/doc/38-aot.md}{AOT}和\href{https://github.com/rabbitc999/ChaLi/blob/master/doc/library/README.md}{标准库}。
